\documentclass[11pt,letterpaper]{article}
\usepackage[T1]{fontenc}
\usepackage{lmodern}
\usepackage[margin=1in,headheight=14pt]{geometry}
\usepackage{amsmath,amssymb,amsthm,mathtools}
\usepackage{booktabs,array,longtable}
\usepackage{xcolor}
\usepackage{microtype}
\usepackage{fancyhdr}
\usepackage{enumitem}
\usepackage{hyperref}
\definecolor{ink}{HTML}{17354C}
\definecolor{pale}{HTML}{F2F5F7}
\hypersetup{colorlinks=true,linkcolor=ink,citecolor=ink,urlcolor=ink,
 pdftitle={IE-27: Recovered analysis and certified finite-stage results},
 pdfauthor={OpenAI ChatGPT}}
\pagestyle{fancy}
\fancyhf{}
\lhead{\small\textcolor{ink}{IE-27 / Recovery edition}}
\rhead{\small 12 September 2026}
\cfoot{\small\thepage}
\setlength{\parindent}{0pt}
\setlength{\parskip}{6pt}
\setlength{\emergencystretch}{2em}
\clubpenalty=10000
\widowpenalty=10000
\newtheorem{theorem}{Theorem}[section]
\newtheorem{lemma}[theorem]{Lemma}
\newtheorem{proposition}[theorem]{Proposition}
\newtheorem{corollary}[theorem]{Corollary}
\theoremstyle{remark}
\newtheorem{remark}[theorem]{Remark}
\newcommand{\R}{\mathbb R}
\newcommand{\C}{\mathbb C}
\newcommand{\Q}{\mathbb Q}
\newcommand{\Id}{I}
\newcommand{\spec}{\sigma}
\newcommand{\norm}[1]{\left\|#1\right\|_2}
\newcommand{\code}[1]{\texttt{\detokenize{#1}}}
\begin{document}

{\small\sffamily\bfseries\textcolor{ink}{MATHEMATICAL RECOVERY REPORT}}

{\LARGE\bfseries\textcolor{ink}{IE-27: Recovered analysis and\\[3pt]certified finite-stage results}}

{\large The spectral disk for a Radau stage preconditioner}

12 September 2026 \hfill Reconstruction of an interrupted investigation

\vspace{5pt}
\fcolorbox{ink}{pale}{\parbox{0.945\linewidth}{
\textbf{Status. This is not a complete solution of IE-27.}
The arbitrary-stage disk inequality has not been established in this work.
The archive contains proved structural statements, exact low-stage arguments,
and 64 freshly verified finite-stage certificates. The covered stages are
\(q=2,3,\ldots,64\), together with \(q=80\). Each certificate covers
\emph{every} positive shift, not a sampled grid.
}}

\begin{abstract}
We reconstruct the mathematical content recoverable from an interrupted
investigation. The original files were absent from the recovery workspace;
the present files and certificate data have been rebuilt and rechecked.
We give a self-contained differentiation formula for the inverse Radau
collocation matrix, prove positivity of every pivot in the specified
factorization, reduce the spatial pencil exactly to stage matrices, and prove
a common-energy sufficient condition for the requested disk. We supply a
closed-form two-stage bound and an explicit three-stage proof in
\(\Q(\sqrt6)\). Integer-interval verification then establishes the disk for
64 specified stage numbers. The proof checker requires only Python's standard
library. A counterexample to a tempting Euclidean resolvent estimate explains
why the elementary norm argument does not solve the general problem. The
remaining all-stage requirement is stated explicitly.
\end{abstract}

\section{Target, notation, and scope}

We use the normalization and symmetric spatial setting of the IE-27 problem
page~\cite{problem}. Let \(A_q\) be the \(q\)-stage Radau IIA collocation matrix,
with increasing nodes \(0<c_1<\cdots<c_q=1\), and write
\begin{equation}\label{eq:factor}
 D_q=A_q^{-1}=L_qU_q,\qquad N_q=U_q-\Id_q,\qquad B_q=L_q^{-1}.
\end{equation}
Here there is no pivoting, \(L_q\) is lower triangular, and \(U_q\) has unit
diagonal. The stipulated radius is \(r_q=\norm{N_q}<1\).
For symmetric positive-definite \(M,K\in\R^{n\times n}\) and \(\tau>0\), define
\begin{equation}\label{eq:pencil}
 \mathcal A=D_q\otimes M+\tau\Id_q\otimes K,
 \qquad \mathcal P_L=L_q\otimes M+\tau\Id_q\otimes K.
\end{equation}
The target is \(|\nu-1|\le r_q\) for every generalized eigenvalue of
\(\mathcal A x=\nu\mathcal P_Lx\), uniformly over the admissible data and stage
number. Related formulations and low-order analyses are in
\cite{dravi,out,axel}. We make no novelty claim for the two- or three-stage
cases.

It is essential not to confuse four different assertions: a spectral-radius
bound, an operator-norm bound, a finite set of stage counts, and an arbitrary
stage count. The results below prove the requested spectral statement at the
listed stages and do not replace its all-stage quantifier by numerical evidence.
Also, establishing \(r_q<1\) for every \(q\) is not the remaining target: IE-27
already assumes that inequality for the stage number under consideration.

\section{An exact construction of the stage matrices}

Write \(\ell_j\) for the Lagrange cardinal polynomials at the nodes. By definition,
\begin{equation}\label{eq:A}
 (A_q)_{ij}=\int_0^{c_i}\ell_j(t)\,dt.
\end{equation}
The interior Radau nodes are the roots of
\begin{equation}\label{eq:J}
 J_q(t)=\sum_{k=0}^{q-1}(-1)^{q-1-k}
 \binom{q-1}{k}\binom{q+k}{k}t^k
       =P_{q-1}^{(1,0)}(2t-1).
\end{equation}
In fact, the usual Legendre description and this one are related by
\(P_q(2t-1)-P_{q-1}(2t-1)=2(t-1)J_q(t)\). This identity follows by comparing
the coefficient formulas, or by the Legendre recurrence. Formula~\eqref{eq:J}
is the integer polynomial used to certify every node in the supplied data.

\begin{lemma}[Inverse as a differentiation matrix]\label{lem:diff}
Let \(R(t)=\prod_{j=1}^q(t-c_j)\) and
\(v_i=c_iR'(c_i)=c_i\prod_{j\ne i}(c_i-c_j)\). Then \(A_q\) is invertible, and
\begin{equation}\label{eq:Dentries}
 (D_q)_{ij}=\frac{v_i}{v_j(c_i-c_j)}\quad(i\ne j),\qquad
 (D_q)_{ii}=\begin{cases}
 (2c_i)^{-1},&i<q,\\
 (q^2+1)/2,&i=q.
 \end{cases}
\end{equation}
\end{lemma}
\begin{proof}
For \(p\) of degree at most \(q-1\), put \(f(t)=\int_0^tp(s)\,ds\).
The vector \((f(c_i))_i\) is \(A_q(p(c_i))_i\). If all \(f(c_i)\) vanish,
then \(f\), a polynomial of degree at most \(q\), has the \(q+1\) distinct
zeros \(0,c_1,\ldots,c_q\). Thus \(f=p=0\), proving invertibility.

Interpolation on the space of polynomials of degree at most \(q\) vanishing
at zero has cardinal basis
\[
 \psi_j(t)=\frac{tR(t)}{c_jR'(c_j)(t-c_j)}.
\]
Differentiation gives \((D_q)_{ij}=\psi_j'(c_i)\). For \(i\ne j\) this is the
first formula in~\eqref{eq:Dentries}; for \(i=j\) it gives
\begin{equation}\label{eq:diag-general}
 (D_q)_{ii}=\frac1{c_i}+\frac{R''(c_i)}{2R'(c_i)}.
\end{equation}

For completeness, the coefficients in~\eqref{eq:J} directly satisfy
\begin{equation}\label{eq:JacobiODE}
 t(1-t)J_q''+(1-3t)J_q'+(q^2-1)J_q=0.
\end{equation}
Indeed, their successive ratio is
\(- (q-1-k)(q+k+1)/(k+1)^2\), which is exactly the coefficient recurrence
in this differential identity.
As \(R\) is a nonzero constant times \((t-1)J_q\), an interior root satisfies
\[
 \frac{R''(c_i)}{2R'(c_i)}
 =\frac1{c_i-1}+\frac{J_q''(c_i)}{2J_q'(c_i)}
 =-\frac1{2c_i}.
\]
At the endpoint, evaluating~\eqref{eq:JacobiODE} at \(t=1\) gives
\(J_q'(1)/J_q(1)=(q^2-1)/2\). Substituting these expressions in
\eqref{eq:diag-general} proves both diagonal formulas.
\end{proof}

\begin{theorem}[Positive pivots for every stage number]\label{thm:pivots}
For every \(q\ge2\), the factorization~\eqref{eq:factor} exists uniquely with
the prescribed normalization, and every diagonal entry of \(L_q\) is positive.
Consequently \(L_q+\mu\Id_q\), and hence \(\Id_q+\mu B_q\), are invertible
for all \(\mu\ge0\).
\end{theorem}
\begin{proof}
Put \(V=\operatorname{diag}(v_1,\ldots,v_q)\) and \(C=V^{-1}D_qV\).
Lemma~\ref{lem:diff} shows that \(C_{ij}=1/(c_i-c_j)\) off the diagonal, so
\[
 C+C^T=2\operatorname{diag}\left(\frac1{2c_1},\ldots,
       \frac1{2c_{q-1}},\frac{q^2+1}{2}\right)\succ0.
\]
Every leading principal submatrix of \(C\) also has a positive-definite
symmetric part. If it has an eigenpair \((\lambda,z)\), then
\(\operatorname{Re}\lambda=\operatorname{Re}(z^*Cz)/(z^*z)>0\).
Its real eigenvalues are positive and its nonreal eigenvalues occur in
conjugate pairs with positive products. Its determinant is therefore positive.
The corresponding leading principal submatrix of \(D_q\) is diagonally
similar to it, so its determinant is positive as well.

Writing these leading determinants as \(\delta_j>0\), with \(\delta_0=1\),
Gaussian elimination without pivoting has pivots
\(\delta_j/\delta_{j-1}>0\). Placing the pivots in the lower factor gives
precisely the normalization in~\eqref{eq:factor}. Uniqueness is the usual
uniqueness of triangular factorization with a prescribed unit diagonal.
The final assertion follows from the positive diagonal entries of the
triangular matrix \(L_q+\mu\Id_q\).
\end{proof}

Diagonal similarity in this proof does \emph{not} preserve the Euclidean norm
of \(N_q\). Thus positivity of the symmetric part of \(C\) cannot simply be
substituted into a proof with radius \(\norm{N_q}\).

\section{Exact removal of the spatial dimension}

\begin{theorem}[Stage reduction]\label{thm:reduction}
For the pencil~\eqref{eq:pencil}, its generalized spectrum is the union,
with algebraic multiplicities, of the spectra of
\begin{equation}\label{eq:X}
 X_{q,\mu}=\Id_q+Y_{q,\mu},\qquad
 Y_{q,\mu}=(\Id_q+\mu B_q)^{-1}N_q,
\end{equation}
where \(\mu\) ranges over the eigenvalues of \(\tau M^{-1}K\). These shifts
are strictly positive.
\end{theorem}
\begin{proof}
Since \(M\) and \(K\) are symmetric positive definite, the symmetric matrix
\(M^{-1/2}KM^{-1/2}\) has an orthogonal diagonalization
\(Q\operatorname{diag}(\kappa_1,\ldots,\kappa_n)Q^T\), where \(\kappa_j>0\).
For \(S=M^{-1/2}Q\) we have
\(S^TMS=\Id_n\) and \(S^TKS=\operatorname{diag}(\kappa_j)\).
Congruence of both members of the pencil by \(\Id_q\otimes S\) preserves
its generalized eigenvalues. A simultaneous permutation then separates it
into the \(n\) pencils
\[
 D_q+\tau\kappa_j\Id_q\quad\hbox{and}\quad L_q+\tau\kappa_j\Id_q.
\]
The latter matrices are invertible by Theorem~\ref{thm:pivots}. Using
\(D_q=L_q(\Id_q+N_q)\), their preconditioned matrices are
\[
 (L_q+\mu\Id_q)^{-1}(D_q+\mu\Id_q)
 =\Id_q+(L_q+\mu\Id_q)^{-1}L_qN_q
 =\Id_q+(\Id_q+\mu B_q)^{-1}N_q.
\]
This proves the assertion including multiplicities.
\end{proof}

This reduction is consistent with the spectral approaches in
\cite{dravi,out}; it is supplied here to make the certification implication
self-contained. In particular, a theorem for all \(\mu>0\) at a fixed stage
immediately applies to every allowed spatial problem and time step. Since
\(N_q\) is strictly upper triangular, \(N_qe_1=0\), so every stage block has
an eigenvalue one. Therefore at least \(n\) eigenvalues of the original
preconditioned system equal one.

\section{The common-energy certificate}

\begin{theorem}[A sufficient metric condition]\label{thm:energy}
Let \(B,N\in\R^{q\times q}\), let \(H=H^T\succ0\), and let \(t\ge0\).
Assume
\begin{equation}\label{eq:LMIs}
 HB+B^TH\succeq0,\qquad tH-N^THN\succeq0.
\end{equation}
Then, for every \(\mu\ge0\), \(\Id_q+\mu B\) is invertible and
\begin{equation}\label{eq:energy-result}
 \rho\big((\Id_q+\mu B)^{-1}N\big)\le\sqrt t.
\end{equation}
\end{theorem}
\begin{proof}
With \(\|z\|_H^2=z^*Hz\), expansion gives
\[
 \|(\Id_q+\mu B)z\|_H^2
 =\|z\|_H^2+\mu z^*(HB+B^TH)z+\mu^2\|Bz\|_H^2
 \ge\|z\|_H^2.
\]
This proves invertibility and that the inverse is a contraction in the
\(H\)-norm. More directly, if
\((\Id_q+\mu B)^{-1}Nz=\lambda z\) with \(z\ne0\), then
\[
 |\lambda|^2\|z\|_H^2
 \le |\lambda|^2\|(\Id_q+\mu B)z\|_H^2
 =\|Nz\|_H^2\le t\|z\|_H^2.
\]
Dividing by the positive quantity \(\|z\|_H^2\) proves the bound.
\end{proof}

\begin{corollary}[A completely checkable sufficient certificate]\label{cor:cert}
For the exact \(B_q,N_q\), suppose rational data \(H=H^T\), \(t\), and a
nonzero rational vector \(v\) satisfy
\begin{align}\label{eq:cert}
&H\succ0,\quad HB_q+B_q^TH\succ0,\quad
 tH-N_q^THN_q\succ0,\quad \Id_q-N_q^TN_q\succ0,\\
&0<t<1,\qquad \|N_qv\|_2^2-t\|v\|_2^2>0.\nonumber
\end{align}
Then for all \(\mu>0\),
\begin{equation}\label{eq:chain}
 \rho(Y_{q,\mu})\le\sqrt t<\norm{N_q}<1.
\end{equation}
Thus the IE-27 disk assertion holds for this stage number and all admissible
spatial matrices and time steps.
\end{corollary}
\begin{proof}
The first three matrix inequalities give Theorem~\ref{thm:energy}.
The Rayleigh quotient of \(N_q^TN_q\) at \(v\) gives \(t<\norm{N_q}^2\).
The fourth matrix inequality gives \(\norm{N_q}<1\).
Theorem~\ref{thm:reduction} transfers the result to the original pencil.
\end{proof}

The strict inequalities are useful for robust interval verification. They are
stronger than necessary for the spectral conclusion. In particular, this is a
\emph{sufficient} certificate condition, not a claimed necessary characterization
of the conjecture.

\section{Exact low-stage arguments}

\subsection{Two stages: a closed-form sharp bound}

For \(q=2\), direct integration and factorization give
\[
 A_2=\begin{pmatrix}5/12&-1/12\\3/4&1/4\end{pmatrix},\qquad
 L_2=\begin{pmatrix}3/2&0\\-9/2&4\end{pmatrix},\quad
 N_2=\begin{pmatrix}0&1/3\\0&0\end{pmatrix},\quad
 B_2=\begin{pmatrix}2/3&0\\3/4&1/4\end{pmatrix}.
\]
It follows that
\[
 Y_{2,\mu}=\begin{pmatrix}
 0&\dfrac1{2\mu+3}\\[5pt]
 0&-\dfrac{3\mu}{(\mu+4)(2\mu+3)}
 \end{pmatrix}.
\]
The eigenvalues are zero and the lower-right entry. Differentiating its
absolute value shows that the maximum over \(\mu>0\) occurs at
\(\mu=\sqrt6\). Hence
\begin{equation}\label{eq:q2}
 \sup_{\mu>0}\rho(Y_{2,\mu})
 =\frac3{11+4\sqrt6}=0.1442449234\ldots <\frac13=\norm{N_2}.
\end{equation}
This supplies a direct proof independent of the certificate generator.
The bundled two-stage certificate uses the looser radius \(\sqrt{21/200}\)
so that the same checker interface covers every included stage.

\subsection{Three stages: an explicit rational metric}

Put \(s=\sqrt6\). The nodes are \((4-s)/10,(4+s)/10,1\), and the exact factors
from Lemma~\ref{lem:diff} yield
\begin{align*}
 B_3&=\begin{pmatrix}
 \frac45-\frac{s}{5}&0&0\\[3pt]
 \frac9{50}+\frac{29s}{200}&\frac3{10}+\frac{3s}{40}&0\\[3pt]
 \frac49-\frac{s}{36}&\frac49+\frac{s}{36}&\frac19
 \end{pmatrix},\\
 N_3&=\begin{pmatrix}
 0&-\frac{53}{25}+\frac{76s}{75}&\frac{16}{25}-\frac{22s}{75}\\[3pt]
 0&0&\frac2{25}+\frac{3s}{25}\\[3pt]
 0&0&0
 \end{pmatrix}.
\end{align*}
Choose the rational data
\begin{equation}\label{eq:q3-H}
 H=\begin{pmatrix}1&0&0\\0&6/5&-7/40\\0&-7/40&23/20\end{pmatrix},
 \qquad t=\frac4{25},\qquad v=(0,1,-1)^T.
\end{equation}
Define \(G=HB_3+B_3^TH\), \(F=tH-N_3^THN_3\), and
\(J=\Id_3-N_3^TN_3\). Table~\ref{tab:minors} gives each leading principal
minor as \(a+bs\). All are strictly positive. The final column is obtained
using only the rational enclosure
\begin{equation}\label{eq:sqrt6}
 \frac{2449489}{10^6}<\sqrt6<\frac{2449490}{10^6},
\end{equation}
whose validity follows by squaring the two positive endpoints.
Thus Sylvester's criterion proves positivity of all four matrices.

\begin{table}[htbp]
\centering
\small
\renewcommand{\arraystretch}{1.6}
\begin{tabular}{@{}ccrrr@{}}
\toprule
Matrix & Minor & $a$ & $b$ & Certified lower bound\\
\midrule
$H$ & 1 & $1$ & $0$ & 1.00000000 \\
$H$ & 2 & $\frac{6}{5}$ & $0$ & 1.20000000 \\
$H$ & 3 & $\frac{2159}{1600}$ & $0$ & 1.34937500 \\
$G$ & 1 & $\frac{8}{5}$ & $-\frac{2}{5}$ & 0.62020400 \\
$G$ & 2 & $\frac{183637549}{648000000}$ & $-\frac{112529}{40500000}$ & 0.27658538 \\
$G$ & 3 & $-\frac{6814458569}{25920000000}$ & $\frac{5594378243}{51840000000}$ & 0.00143616 \\
$F$ & 1 & $\frac{4}{25}$ & $0$ & 0.16000000 \\
$F$ & 2 & $-\frac{78476}{46875}$ & $\frac{32224}{46875}$ & 0.00973511 \\
$F$ & 3 & $-\frac{6359599}{29296875}$ & $\frac{2597464}{29296875}$ & 0.00009763 \\
$J$ & 1 & $1$ & $0$ & 1.00000000 \\
$J$ & 2 & $-\frac{18104}{1875}$ & $\frac{8056}{1875}$ & 0.86884447 \\
$J$ & 3 & $-\frac{19088}{1875}$ & $\frac{1672}{375}$ & 0.74118828 \\
\bottomrule
\end{tabular}

\caption{Exact three-stage leading principal minors, written \(a+b\sqrt6\).
The displayed decimal lower bounds are rounded downward; every bound is positive.}
\label{tab:minors}
\end{table}

For the vector in~\eqref{eq:q3-H}, direct calculation gives
\begin{equation}\label{eq:q3gap}
 \|N_3v\|_2^2-\frac4{25}\|v\|_2^2
 =\frac{6613}{375}-\frac{4496\sqrt6}{625}>0,
\end{equation}
again certified by~\eqref{eq:sqrt6}. Corollary~\ref{cor:cert} therefore proves
\begin{equation}\label{eq:q3result}
 \boxed{\displaystyle
 \rho(Y_{3,\mu})\le\frac25<\norm{N_3}<1\quad\text{for every }\mu>0.}
\end{equation}
This radius is not claimed to be optimal. The independent program
\code{code/exact_q3.py} performs the matrix identity \(B_3D_3=\Id_3+N_3\),
all 12 determinant calculations, and the sign tests entirely in
\(\Q(\sqrt6)\), using Python integer fractions. It uses neither an eigensolver
nor a symbolic-algebra package.

\section{Fresh finite-stage certification}

\begin{theorem}[Verified finite set]\label{thm:finite}
Let
\[
 \mathcal Q=\{2,3,\ldots,64\}\cup\{80\}.
\]
For every \(q\in\mathcal Q\), the certificate included in
\code{certificates/} satisfies~\eqref{eq:cert}. In particular,
\(\rho(Y_{q,\mu})\le\sqrt{t_q}<\norm{N_q}<1\) for every \(\mu>0\).
Consequently the stated IE-27 disk holds at every one of these 64 stages,
for every admissible \(M,K,n,\tau\).
\end{theorem}
\begin{proof}[Proof by integer-interval certificates]
The supplied data and the accepted execution record in
\code{results/verified.json} establish~\eqref{eq:cert} as follows.

For each \(q\), the certificate specifies \(q-1\) disjoint, increasing dyadic
intervals inside \((0,1)\). The checker evaluates the integer polynomial
\eqref{eq:J} at their endpoints using exact integer Horner evaluation and
checks strict sign changes. An interval reduced to a point is allowed only
when exact evaluation is zero. A polynomial of degree \(q-1\) with a root
in each of these \(q-1\) disjoint intervals has precisely one simple root
in each, and no remaining roots. Appending the exact node one identifies
the Radau nodes, not rounded substitutes.

Every subsequent scalar is enclosed by an interval
\([a/2^{1024},b/2^{1024}]\) with integer endpoints. Addition, multiplication,
squaring, and division use outward rounding by exact integer floor and ceiling.
Division is refused whenever the denominator interval contains zero. Substituting
the node enclosures in~\eqref{eq:Dentries}, the checker performs the normalized
triangular factorization and forward substitution to enclose the exact
\(B_q,N_q\). Each lower-factor pivot is also checked to have a strictly
positive lower bound.

The supplied rational \(H,t,v\) are read exactly. Their dimensions, exact
symmetry of \(H\), and \(0<t<1\) are checked. The four symmetric exact
matrices in~\eqref{eq:cert} are then enclosed and tested by interval
\(LDL^T\) elimination. Every exact pivot is contained in its computed
interval; strict positivity of all the lower endpoints certifies positive
definiteness by congruence. Finally the checker computes an interval for
\(\|N_qv\|_2^2-t\|v\|_2^2\) and requires its lower endpoint to be positive,
with \(v\ne0\).

All 64 included files pass these tests in the supplied fresh execution.
No floating-point eigenvalue, optimizer output, or approximate root is used
as a proof decision. Corollary~\ref{cor:cert} supplies the conclusion for the
whole half-line of shifts.
\end{proof}

\begin{table}[htbp]
\centering
\small
\renewcommand{\arraystretch}{1.6}
\begin{tabular}{@{}crrr@{}}
\toprule
$q$ & Exact $t_q$ & $\sqrt{t_q}$ (approx.) & Lower bound on energy pivot\\
\midrule
2 & $\frac{21}{200}$ & 0.324037035 & 0.039531249999 \\
3 & $\frac{4}{25}$ & 0.400000000 & 0.005192775643 \\
4 & $\frac{18235523}{80000000}$ & 0.477434852 & 0.000631320126 \\
8 & $\frac{3842650859}{10000000000}$ & 0.619891189 & 0.000154975869 \\
16 & $\frac{263342139}{500000000}$ & 0.725730169 & 0.000237051813 \\
32 & $\frac{6442271059}{10000000000}$ & 0.802637593 & 0.000160167625 \\
48 & $\frac{3505623359}{5000000000}$ & 0.837331877 & 0.000287967193 \\
64 & $\frac{7365589189}{10000000000}$ & 0.858230108 & 0.000042466148 \\
80 & $\frac{1522846637}{2000000000}$ & 0.872595736 & 0.000170002136 \\
\bottomrule
\end{tabular}

\caption{Representative accepted certificates. The last column bounds below the
smallest \(LDL^T\) pivot of \(HB_q+B_q^TH\), not its smallest eigenvalue.
The middle decimal column is illustrative only; the proof uses the exact rational \(t_q\).}
\label{tab:certs}
\end{table}

Table~\ref{tab:certs} is a selection, not the full verification set.
Complete machine-readable results appear in the \code{results/} directory,
in \code{certificates.csv} and \code{verified.json}. They include all 64
cases and all strict margins. The gaps between stages 64 and 80 are intentional:
\emph{this archive does not certify stages 65--79}.

\subsection{How candidates were found}

The generator computes approximate stage matrices in a shifted Legendre basis,
then seeks a symmetric tridiagonal metric. For \(q\ge4\), it takes a rational
\(t_q\) near \((0.999\,\norm{N_q})^2\), where the norm here is an untrusted
numerical estimate, and minimizes squared negative eigenvalues of
\[
 H-\Id_q,\qquad q(B_q^TH+HB_q)-10^{-3}\Id_q,\qquad
 t_qH-N_q^THN_q-10^{-3}\Id_q.
\]
Multiplication of the energy matrix by the positive scalar \(q\) changes no
positivity condition. Candidate entries of \(H\) and a numerical right
singular vector of \(N_q\) are rounded to rational numbers with denominator
\(10^8\). Integer bisection then supplies the root brackets. The exact
low-stage choices override the optimizer at \(q=2,3\).

These are only proposal procedures. Their floating-point computations and
termination flags are not trusted. The independent checker establishes the
inequalities after rounding, or rejects the candidate. A failed search or a
failed interval test would not by itself refute the conjecture.

\section{Why the elementary resolvent proof fails}

A tempting argument is
\[
 \rho((\Id_q+\mu B_q)^{-1}N_q)
 \le\norm{(\Id_q+\mu B_q)^{-1}}\,\norm{N_q}
 \stackrel{?}{\le}\norm{N_q}.
\]
The second step would follow from Euclidean accretivity of \(B_q\), but that
property already fails at three stages.

\begin{proposition}[Failure of Euclidean resolvent contractivity]\label{prop:failure}
For \(q=3\), \(\norm{(\Id_3+\mu B_3)^{-1}}>1\) for all sufficiently small
positive \(\mu\).
\end{proposition}
\begin{proof}
For \(w=(0,1,-2)^T\), the explicit matrix above gives
\[
 \alpha=w^TB_3w=\frac{-52+7\sqrt6}{360}<0.
\]
Set \(\beta=\|B_3w\|_2^2>0\). Whenever \(0<\mu<-2\alpha/\beta\),
\[
 \|(\Id_3+\mu B_3)w\|_2^2
 =\|w\|_2^2+2\mu\alpha+\mu^2\beta<\|w\|_2^2.
\]
The matrix is invertible, so applying its inverse to the nonzero vector
\((\Id_3+\mu B_3)w\) proves the asserted strict norm inequality.
\end{proof}

This is a counterexample to the \emph{shortcut}, not to IE-27.
The three-stage spectral disk has just been proved by a different norm.
Likewise, replacing a matrix by a diagonally similar one can preserve its
spectrum while changing the particular Euclidean radius in the conjecture.
Such a change of metric must be justified, not silently identified with
an orthogonal transformation.

\section{The exact remaining gap}

To turn the common-energy route into an all-stage solution, one possible
sufficient result would be:
\begin{quote}
For every \(q\ge2\) satisfying \(\norm{N_q}<1\), there is a symmetric
\(H_q\succ0\) such that
\[
 H_qB_q+B_q^TH_q\succeq0,\qquad
 N_q^TH_qN_q\preceq\norm{N_q}^2H_q.
\]
\end{quote}
Theorem~\ref{thm:energy} would then give the desired result immediately.
\textbf{No proof of this assertion is supplied here.} In particular, the
tridiagonal optimization successes do not furnish an analytic formula or a
verified existence argument for all stage numbers.

This sufficient assertion need not be necessary. A direct proof bounding the
roots of the generalized characteristic polynomial, or a counterexample at an
admissible stage and positive shift, could settle the target without a common
metric. Neither has been established by this recovery. Repeated finite
verification does not eliminate the universal quantifier, and the positive-pivot
theorem alone does not imply the disk radius.

The precise conclusion of the archive is therefore Theorem~\ref{thm:finite},
not a declaration that IE-27 has been fully solved.

\section{Reproducibility and trust boundary}

The main commands, run from the extracted archive directory, are
\begin{verbatim}
python code/run_checks.py
python code/exact_q3.py
python code/verify.py certificates/q080.json
\end{verbatim}
The first command executes the 12 regression tests, the independent three-stage
proof, and every certificate. It writes new results to \code{recheck/}, leaving
the delivered baseline logs unchanged. Verification needs Python 3.10 or newer
and no third-party libraries. The optional generator uses NumPy and SciPy;
the generation environment is recorded in \code{results/environment.txt}.

\code{code/test_verify.py} tests exact fractional enclosure on random rational
inputs, interval extrema and reciprocals, integer polynomial evaluation,
positive and nonpositive \(LDL^T\) examples, exact two-stage entries, and rejection
of deliberately corrupted certificates. All 12 tests passed during recovery.
The tests are useful safeguards, not a substitute for the mathematical argument
or a formal proof of the checker implementation.

The trusted components are the formulas proved in this report, the transparent
interval checker and exact-field checker, and the usual correctness of Python
integer arithmetic and execution. This is a computer-assisted proof with
inspectable code, not a proof-assistant formalization. The files
\code{MANIFEST.sha256} and \code{RECOVERY.md} record integrity and provenance.
No claim of peer review, exhaustive literature search, or required work duration
is made. The primary-source references identify the target and its context.

\begingroup
\small
\begin{thebibliography}{9}
\setlength{\itemsep}{3pt}
\bibitem{problem}
OpenProblemsInNLA, \emph{IE-27 --- Spectral disk for a Radau stage preconditioner},
problem README read directly on 12 September 2026.
\href{https://raw.githubusercontent.com/ajt60gaibb/OpenProblemsInNLA/main/linear-systems-and-elimination/IE-27/README.md}{Repository problem statement (raw README)}.

\bibitem{dravi}
I.~Dravins, S.~Serra-Capizzano, and M.~Neytcheva,
\emph{Fine spectral analysis of preconditioned matrices and matrix-sequences
arising from stage-parallel implicit Runge-Kutta methods of arbitrarily high
order}, arXiv:2302.04657v2, 24 February 2023, especially Sections 2 and 4.
\url{https://arxiv.org/html/2302.04657v2}.

\bibitem{out}
M.~Outrata, \emph{On the recent advances of spectral analysis for systems
arising from fully-implicit RK methods}, arXiv:2510.21241v1, 24 October 2025,
especially Sections 2.1--2.3.
\url{https://arxiv.org/html/2510.21241v1}.

\bibitem{axel}
O.~Axelsson, I.~Dravins, and M.~Neytcheva,
\emph{Stage-parallel preconditioners for implicit Runge-Kutta methods of
arbitrarily high order, linear problems}, Numerical Linear Algebra with
Applications 31(1), e2532 (2024). DOI: 10.1002/nla.2532.
Bibliographic details and conjecture attribution from the problem page;
a fresh full-text inspection is not claimed.
\url{https://doi.org/10.1002/nla.2532}.
\end{thebibliography}
\endgroup

\end{document}
